\emph{Setting.} The field content is the pair of Definition \claimref{obj.synthesis.gtg-flat-spacetime-gauge-fields}: a position-gauge field $\hb$ and a rotation-gauge field $\Omega$, with the rotation gauge $\hb^{-1}(a) \mapsto U\hb^{-1}(a)\tilde{U}$, $\Omega(a) \mapsto U\Omega(a)\tilde{U} - 2(a \cdot \nabla U)\tilde{U}$ and the displacement gauge acting on $\hb^{-1}$ and $\Omega$ as one-forms (the input slot transforms by the differential $\underline{f}$, the values are unchanged). Rotors of $W$ commute with $\nO$ and $\nInf$, since these are orthogonal to $W$; so the rotation gauge acts on $\Cl(2,4)$, fixing the plane of $\nO$ and $\nInf$. All statements below are pointwise linear algebra in $\Cl(2,4)$, apart from the variational step in (i). The map $\Phi_{\rho}(\hb, \Omega) = \Omega - u_{\rho} \wedge \hb^{-1}$ of the statement is the field space throughout; clause (ii) shows it is exactly the gauge fields valued in the stabiliser of $I_{\rho}$.

\emph{(i).} Fix $\rho \neq 0$, $I = I_{\rho}$, $V = -u_{\rho}/(2\rho)$ (so $Q(V) = -1/(2\rho) = -3/\Lambda$ and $V \cdot I = 0$). On the field space $\Phi_{\rho}$, $D_{a}I = A(a) \times I = 0$, and $D_{a}V = A(a) \times V$ with $(u_{\rho} \wedge c) \times u_{\rho} = -Q(u_{\rho})c = 2\rho c$ (checked), so $D_{a}V$ is a nonzero multiple of $E(a)$, the soldering form. Vary $\Omega$ by $\delta\Omega$ with compact support: $\delta F = D\delta\Omega$ and $\delta S \propto 2\int\epsilon(I, V, D\delta\Omega, F)$. Since $\epsilon$ is invariant under $\so(2,4)$, $d\,\epsilon(I, V, \delta\Omega, F) = \epsilon(DI, V, \delta\Omega, F) + \epsilon(I, DV, \delta\Omega, F) + \epsilon(I, V, D\delta\Omega, F) \pm \epsilon(I, V, \delta\Omega, DF)$ with the graded signs of forms; $DI = 0$ and $DF = 0$ (Bianchi), so up to a boundary term the $\Omega$-equation is $\epsilon(I, E, \delta\Omega, F) = 0$ for all $\Lambda^{2}W$-valued $\delta\Omega$. The check shows $\epsilon(I_{\rho}, w, B, C) = 0$ for all $w \in W$, $B, C \in \Lambda^{2}W$ (the six vectors lie in the five-dimensional span of $I_{\rho}$ and $W$), so only the block $-u_{\rho} \wedge T$ of $F$ enters, and it computes the resulting linear map from the twenty-four torsion components to the twenty-four components of the $\Omega$-equation (contracting with the background Levi-Civita symbol) at the identity frame and at a generic rational frame: rank $24$ in both, for $\rho = \tfrac12, -\tfrac12, \tfrac{3}{14}$. Thus without matter the $\Omega$-equation is $T = 0$, the first Einstein-Cartan equation of Fact \claimref{dark-energy.imported-macdowell-mansouri-theorem-stelle-west}, and it holds identically on the subfamily $\Omega = \Omega(\hb)$. On that subfamily the action is a functional of $\hb$ alone, and by the chain rule its variation is the frame variation at fixed $\Omega$ plus the $\Omega$-equation contracted with $\delta\Omega(\hb)$; the second term vanishes identically there (the $\Omega$-equation is zero pointwise, so the derivatives of $\delta\hb$ inside $\delta\Omega(\hb)$ do not matter). So its Euler-Lagrange equation is the frame equation on torsion-free fields, which is $2(\mathcal{G}(a) + \Lambda a)I = 0$ by Theorem \claimref{synthesis.macdowell-mansouri-principle-fixes-the-trace-free-part} (rank $10$ on the algebraic curvature tensors with kernel the Weyl part, Theorem \claimref{synthesis.curvature-only-tractor-actions-are-macdowell-mansouri}(iv)), that is $\mathcal{G}(a) + \Lambda a = 0$ with $\Lambda = 6\rho$. The tractor gauge field of the same torsion-free field is $\Omega(\hb) - \nInf \wedge \hb^{-1} + \nO \wedge \Sch(\hb^{-1})$, so it differs from $\Phi_{\rho}(\hb, \Omega(\hb))$ only in the $\nO$ slot (and, by (ii) here with Theorem \claimref{synthesis.gtg-tractor-gauge-field-reproduces-tractor-connection}(ii), agrees with it exactly where $\Sch = \rho\,\mathrm{id}$), and on it the integrand is $\ip{X}{I}\tfrac14 W_{abcd}W^{abcd}$ with the Bach equations (Theorem \claimref{synthesis.frozen-point-integrand-on-tractor-transport-is-weyl-squared}(iii)). The two one-field spaces have the same rotation-gauge field and the same soldering; the slot alone separates Einstein's equations from Bach's. This is the frozen row against the Schouten row of the slot lattice of Theorem \claimref{synthesis.frozen-point-integrand-gives-einstein-only-with-the-symmetric-slot-frozen}, whose restriction to torsion-free fields on the frozen rows is what (i) justifies. With matter coupled to $\Omega$ the $\Omega$-equation acquires the spin density: independent $\Omega$ then gives torsion proportional to spin (Einstein-Cartan, Fact \claimref{synthesis.imported-einstein-cartan-matter-coupling}), while the subfamily $\Omega = \Omega(\hb)$ keeps the torsion-free coupling. The general identification of first-order and metric variation without spin is standard (cited from memory); the argument above needs only the pointwise computation.

\emph{(ii).} Write a bivector as $B = B_{W} + \nInf \wedge p + \nO \wedge q + \beta\,\nInf \wedge \nO$ with $B_{W} \in \Lambda^{2}W$, $p, q \in W$. For a vector $J = \alpha\nInf + \gamma\nO$ of the plane, $B \times J = B \cdot J$ and $(\nInf \wedge p) \cdot J = \nInf (p \cdot J) - p(\nInf \cdot J) = -\gamma p$, $(\nO \wedge q) \cdot J = -\alpha q$, $(\nInf \wedge \nO) \cdot J = \alpha\nInf - \gamma\nO$, nonzero unless $J = 0$, and $B_{W} \cdot J = 0$. So $B \times I_{\rho} = 0$ ($\alpha = 1$, $\gamma = \rho$) iff $\beta = 0$ and $q = -\rho p$, i.e. $B = B_{W} + \nInf \wedge p - \rho\nO \wedge p = B_{W} + u_{\rho} \wedge p$. The check computes the annihilator of $I_{\rho}$ in the fifteen bivectors exactly (dimension $10$) and shows that $\Lambda^{2}W$ and $u_{\rho} \wedge W$ lie in it with rank $10$, for $\rho \in \{0, \tfrac12, -\tfrac12, \tfrac{3}{14}\}$. A vector annihilated by all of $\Lambda^{2}W$ lies in the plane of $\nO$, $\nInf$, and by the computation above $u_{\rho} \wedge w$ annihilates $\alpha\nInf + \gamma\nO$ iff $\gamma = \rho\alpha$, so the common annihilator is the line of $I_{\rho}$ (the check computes this kernel in all six vectors). The map $\Phi_{\rho}$ is therefore a pointwise linear bijection between $(\hb^{-1}(a), \Omega(a))$ and the values of an $I_{\rho}$-preserving field; its soldering part $p = -\hb^{-1}(a)$ is nondegenerate iff $\det\hb \neq 0$. Equivariance: $U(u_{\rho} \wedge E)\tilde{U} = u_{\rho} \wedge (UE\tilde{U})$, and the inhomogeneous term $-2(a \cdot \nabla U)\tilde{U}$ lies in $\Lambda^{2}W$, where $\Phi_{\rho}$ is the identity; the check confirms $U\Phi_{\rho}\tilde{U} = \Phi_{\rho}(UE\tilde{U}, U\Omega\tilde{U})$ for three rational rotors (two boosts, one rotation) that fix $\nO$, $\nInf$. Field strength: since the commutator with $\Omega(a) \in \Lambda^{2}W$ is a derivation and commutes with $u_{\rho}$, $\Omega(a) \times (u_{\rho} \wedge E(b)) = u_{\rho} \wedge (\Omega(a) \times E(b))$, and for $c, d \in W$, $(u_{\rho} \wedge c) \times (u_{\rho} \wedge d) = -Q(u_{\rho})\,c \wedge d = 2\rho\, c \wedge d$. Substituting $A = \Omega - u_{\rho} \wedge E$ in $F(a \wedge b) = a \cdot \nabla A(b) - b \cdot \nabla A(a) + A(a) \times A(b)$ gives $R(a \wedge b) + 2\rho E(a) \wedge E(b) - u_{\rho} \wedge (D_{a}E(b) - D_{b}E(a))$. The check recomputes this from first jets of $(E, \Omega)$ at a point (only values and first derivatives enter $F$), for the four values of $\rho$, and confirms $F \times I_{\rho} = 0$ with nonzero torsion. With $a \wedge b = \hb(B)$, $E(a) \wedge E(b) = B$ (outermorphism) and $R(\hb(B)) = -\mathcal{R}(B)$, the Lorentz block is $-\mathcal{R}(B) + 2\rho B = -(\mathcal{R}(B) - \tfrac{\Lambda}{3}B)$. At $\rho = 0$ the brackets of the translation block vanish (checked), giving $\mathfrak{iso}(1,3)$.

\emph{Trap check.} Clause (ii) uses only the definition of the gravitational field, whose past cone is empty, and the chart and gauge laws as set up in Theorem \claimref{synthesis.gtg-tractor-gauge-field-reproduces-tractor-connection}, not its clause (ii); no Einstein-equation import is used as a premise. Clause (i) compares field equations on two field spaces given in advance and imposes neither the Einstein condition nor a parallel scale on the normal tractor connection, which enters only as the contrast with the Bach equations. Why the varied field should be $\Phi_{\rho}$ rather than the normal tractor connection is not derived here: the field space is a premise.
